\chapter{SYSTEM IMPLEMENTATION AND TESTING}

\section{Technology Stack \& Project Structure}
The system implementation divides responsibilities between a Python service backend for real-time computer vision and a React-based edge application frontend.
\begin{table}[h]
\centering
\begin{tabular}{p{3.5cm} p{4cm} p{8.5cm}}
\toprule
\textbf{Layer} & \textbf{Technology Used} & \textbf{Role / Context} \\ \midrule
User Interface & React 19, HTML5, Vite & Renders the glassmorphic pages, playground dashboard, and media controls. \\ \midrule
Style Sheet & Vanilla CSS3 & Handles the premium dark-mode layout, glass frosted cards, and animations. \\ \midrule
Web Server & Python 3, Flask & Directs API endpoints, manages webcam streams, and handles asynchronous video uploads. \\ \midrule
Pose Estimation & Google MediaPipe Pose & Locally processes camera frames to locate 33 body joint coordinates. \\ \midrule
Joint Math & NumPy, math libraries & Computes joint flexion-extension angles via 3D vector trigonometry. \\ \midrule
Sequence Classifier & TensorFlow, Keras & Runs 1D-CNN + LSTM inference on 30-frame coordinate vectors. \\ \bottomrule
\end{tabular}
\caption{System Implementation Technology Stack.}
\label{tab:technology_stack}
\end{table}

\section{Module-Wise Implementation Description}

\subsection{Biomechanical Joint Angle Calculations}
The \texttt{BodyPartAngle} class in \texttt{body\_part\_angle.py} computes joint flexion angles. Given three coordinates representing the joints (e.g., Left Shoulder $A$, Left Elbow $B$, and Left Wrist $C$), the angle at the middle joint $B$ is derived using 3D vector trigonometry. 

Let the three joints be represented by vectors:
\[ \vec{A} = (x_A, y_A, z_A), \quad \vec{B} = (x_B, y_B, z_B), \quad \vec{C} = (x_C, y_C, z_C) \]
We compute the vectors connecting elbow to shoulder and wrist:
\[ \vec{u} = \vec{A} - \vec{B} = (x_A - x_B, y_A - y_B, z_A - z_B) \]
\[ \vec{v} = \vec{C} - \vec{B} = (x_C - x_B, y_C - y_B, z_C - z_B) \]
The angle $\theta$ (in radians) is calculated using the dot product:
\[ \cos(\theta) = \frac{\vec{u} \cdot \vec{v}}{\|\vec{u}\| \|\vec{v}\|} \]
\[ \theta = \arccos\left( \frac{u_x v_x + u_y v_y + u_z v_z}{\sqrt{u_x^2 + u_y^2 + u_z^2} \sqrt{v_x^2 + v_y^2 + v_z^2}} \right) \]
We then convert $\theta$ to degrees:
\[ \theta_{\text{deg}} = \theta \times \frac{180}{\pi} \]
To ensure values lie in $[0, 180]$, we handle reflections:
\[ \text{if } \theta_{\text{deg}} > 180^{\circ}, \quad \theta_{\text{deg}} = 360^{\circ} - \theta_{\text{deg}} \]

The implementation of this calculation is shown in Listing 4.1.

\begin{lstlisting}[language=Python, caption=Mathematical Angle Calculation from Joint Vectors, label=lst:angle_calc]
import numpy as np

def calculate_angle(a, b, c):
    a = np.array(a)  # First point
    b = np.array(b)  # Mid point (joint)
    c = np.array(c)  # End point

    radians = np.arctan2(c[1] - b[1], c[0] - b[0]) - \
              np.arctan2(a[1] - b[1], a[0] - b[0])
    angle = np.abs(radians * 180.0 / np.pi)

    if angle > 180.0:
        angle = 360.0 - angle

    return angle
\end{lstlisting}

\subsection{Repetition Counting State Machine}
To count repetitions without double-triggering during shaky movements, the \texttt{TypeOfExercise} class in \texttt{types\_of\_exercise.py} implements threshold-based state machines. Each exercise contains two states:
\begin{itemize}
    \item \textbf{Status = True (DOWN state):} Represents the flexed phase of the movement. For example, in a push-up, the arm angle must drop below $70^{\circ}$. Once triggered, the state switches to False.
    \item \textbf{Status = False (UP state):} Represents the recovery or extension phase. The arm angle must exceed $160^{\circ}$ to reset the state back to True and allow the next repetition to count.
\end{itemize}
To filter out high-frequency noise, a temporal refractory gate checks that at least $0.9$ seconds have elapsed between consecutive repetitions. 

\begin{lstlisting}[language=Python, caption=Squat Repetition State Machine Code, label=lst:squat_state]
def squat(self, counter, status):
    left_leg_angle = self.angle_of_the_right_leg()
    right_leg_angle = self.angle_of_the_left_leg()
    avg_leg_angle = (left_leg_angle + right_leg_angle) // 2

    # Status=True corresponds to ready to trigger
    if status:
        if avg_leg_angle < 70:
            if can_count_rep():  # 0.9s delay check
                counter += 1
                status = False
    else:
        time_diff = time.time() - _last_rep_time["timestamp"]
        if avg_leg_angle > 160 and (time_diff > 0.5):
            status = True

    return [counter, status]
\end{lstlisting}

\subsection{Posture Guard Form Check}
The \texttt{pose\_correction.py} script monitors joints to output visual warnings. Standard checks include:
\begin{itemize}
    \item \textbf{Shoulder Level Guard:} If the absolute vertical difference between the left and right shoulder exceeds $0.03$ normalized units, the prompt ``Straighten Your Shoulders'' is displayed.
    \item \textbf{Hip Sagging Guard:} In planks and push-ups, the difference in height between shoulders and hips is validated; a difference exceeding $0.05$ triggers the warning ``Keep Your Body Straight''.
    \item \textbf{Elbow Flare Check:} In push-ups, checking that elbow coordinates remain close to the torso line ensures proper arm angle.
\end{itemize}

\subsection{Flask MJPEG Streaming \& Sequence Classifier}
The Flask application in \texttt{main.py} runs the video pipeline. It loads the Keras model weights at boot time. For webcam feeds, it processes frames in a generator loop, overlays skeletons and metrics, and returns the multipart HTTP stream.

\begin{lstlisting}[language=Python, caption=Flask Stream Generator Loop, label=lst:flask_stream]
def generate():
    with mp_pose.Pose(
        min_detection_confidence=0.5, 
        min_tracking_confidence=0.5, 
        model_complexity=0
    ) as pose:
        counter = 0
        status = True
        while True:
            ret, frame = cap.read()
            if not ret:
                break
            frame = cv2.resize(frame, (640, 360))
            frame_rgb = cv2.cvtColor(frame, cv2.COLOR_BGR2RGB)
            results = pose.process(frame_rgb)
            
            if results.pose_landmarks:
                landmarks = results.pose_landmarks.landmark
                # Run exercise Rep calculations
                counter, status = (TypeOfExercise(landmarks)
                                   .calculate_exercise(exercise_type, 
                                                       counter, status))
                # Draw skeleton
                mp_drawing.draw_landmarks(
                    frame, results.pose_landmarks, 
                    mp_pose.POSE_CONNECTIONS
                )
            
            # Compress and stream frame
            _, buffer = cv2.imencode(
                '.jpg', frame, 
                [int(cv2.IMWRITE_JPEG_QUALITY), 70]
            )
            yield (b'--frame\r\n'
                   b'Content-Type: image/jpeg\r\n\r\n' 
                   + buffer.tobytes() + b'\r\n')
\end{lstlisting}

\subsection{React UI Client Dashboard}
The React application divides page structures into modular paths. The \texttt{Playground.jsx} page uses the HTML \texttt{<img>} tag to point directly to the backend's stream route (e.g., \url{http://localhost:5000/start-webcam} with query parameter \texttt{?exercise\_type=push-up}), rendering the annotated computer vision output. State values are requested via API polling to update the React target progress indicators and workout HUD statistics in real-time.

\section{System Testing}

\subsection{Testing Objectives}
The testing phase aimed to verify:
\begin{enumerate}
    \item MediaPipe keypoint detection stability across different body distances and camera angles.
    \item Joint angle calculation consistency in static and dynamic postures.
    \item Repetition counter correctness (eliminating false repetitions from non-workout motions).
    \item Deep learning model accuracy in classifying the 13 exercise classes.
    \item Interface synchronization between the Flask streaming thread and React indicators.
\end{enumerate}

\subsection{Test Cases Matrix}
Table 4.4 lists the major test cases executed to validate system modules:
\begin{table}[h]
\centering
\begin{tabular}{|p{1.4cm}|p{2.2cm}|p{4.8cm}|p{4.8cm}|p{1.4cm}|}
\hline
\textbf{Test ID} & \textbf{Component} & \textbf{Input / Action} & \textbf{Expected Result} & \textbf{Status} \\ \hline
TC-01 & Base API & Send GET request to \texttt{/} endpoint. & Returns HTTP 200 with online status message. & Passed \\ \hline
TC-02 & Pose Detection & Webcam feed containing standing person. & MediaPipe localizes all 33 keypoints correctly. & Passed \\ \hline
TC-03 & Angle Calculation & Flex left arm at exactly $90^{\circ}$ check. & Calculated angle is within $\pm 2.5^{\circ}$ accuracy. & Passed \\ \hline
TC-04 & Push-up Rep Count & Flex arm below $70^{\circ}$ and extend past $160^{\circ}$. & Rep counter increments by 1. & Passed \\ \hline
TC-05 & Double Count Filter & Fast arm flex-extend under $0.5$ seconds. & Refractory period filters secondary count. & Passed \\ \hline
TC-06 & Posture Warning & Push-up with arched lower back. & Displays ``WARNING: Wrong Exercise Pose!''. & Passed \\ \hline
TC-07 & Video Analyzer & Drag-and-drop $30$s push-up MP4 file. & Upload succeeds, video streams back with overlays. & Passed \\ \hline
TC-08 & Auto Exercise Detect & Execute bicep curl movements. & Sequence classifier updates active exercise to bicep-curl. & Passed \\ \hline
TC-09 & Goal Completion & Counter reaches selected Target Reps. & Displays goal celebration overlay and completion icon. & Passed \\ \hline
TC-10 & Stream Close & Toggle webcam button to inactive. & Camera feed terminates cleanly, releasing system hardware resources. & Passed \\ \hline
\end{tabular}
\caption{System Representative Test Cases Matrix.}
\label{tab:test_cases}
\end{table}
\newpage
